\documentclass[article]{jss}

%% -- LaTeX packages and custom commands ---------------------------------------
\usepackage{orcidlink,thumbpdf}
\IfFileExists{lmodern.sty}{\usepackage{lmodern}}{}
\usepackage{amsmath,bm,booktabs}

\newcommand{\class}[1]{`\code{#1}'}
\newcommand{\fct}[1]{\code{#1()}}
%% placeholders for the author -- remove before submission
\newcommand{\todo}[1]{\textcolor{red}{\textbf{[#1]}}}

%% -- Article metainformation --------------------------------------------------
\author{Abdul Amir Taima Bandar~\orcidlink{0009-0007-1118-1993}\\University of Kerbala
   \And Ahmed Turki Abd Ali~\orcidlink{0009-0007-5402-8948}\\University of Kerbala}
\Plainauthor{Abdul Amir Taima Bandar, Ahmed Turki Abd Ali}

\title{\pkg{spdfiv}: Defactored Instrumental Variables Estimation and Heterogeneity Diagnostics for Spatial Dynamic Panels with Interactive Effects in \proglang{Python}}
\Plaintitle{spdfiv: Defactored Instrumental Variables Estimation and Heterogeneity Diagnostics for Spatial Dynamic Panels with Interactive Effects in Python}
\Shorttitle{\pkg{spdfiv}: Spatial Dynamic Panels with Interactive Effects in \proglang{Python}}

\Abstract{
Spatial dynamic panel data models with unobserved common factors can be
estimated by the two-stage instrumental variables (2SIV) estimator of Cui,
Sarafidis, and Yamagata (2023) and, with heterogeneous slopes, by the
mean-group IV estimator of Chen, Cui, Sarafidis, and Yamagata (2025). This
article introduces \pkg{spdfiv}, an open-source \proglang{Python}
implementation of both estimators, to our knowledge the first outside
\proglang{Stata}. The package reproduces the published
\proglang{Stata} results of Kripfganz and Sarafidis (2025) in every printed
digit; the comparison uncovered five implementation details that are not
visible from the published equations, which we document. The package also
matches the weights matrix to the panel by unit label, computes impacts
exactly through a Schur decomposition with Delta-method or simulation
inference, provides diagnostics that show which units drive a mean-group
estimate, and offers an experimental treatment of unbalanced panels. Monte
Carlo experiments document each feature. Applied to the bank data of
Kripfganz and Sarafidis (2025), the diagnostics show that extreme bank-level
estimates explain most of the gap between the mean-group and pooled
estimates of some slopes, but not the disappearance of the spatial
spillover.
}

\Keywords{panel data, spatial econometrics, interactive effects, common factors, instrumental variables, mean-group estimation, outliers, \proglang{Python}}
\Plainkeywords{panel data, spatial econometrics, interactive effects, common factors, instrumental variables, mean-group estimation, outliers, Python}

\Address{
  Abdul Amir Taima Bandar (corresponding author), Ahmed Turki Abd Ali\\
  Department of Statistics\\
  College of Administration and Economics\\
  University of Kerbala\\
  Karbala 56001, Iraq\\
  E-mail: \email{Abdulameer.t@uokerbala.edu.iq}, \email{Ahmed.turki@uokerbala.edu.iq}
}

\begin{document}

%% ============================================================================
\section{Introduction} \label{sec:intro}
%% ============================================================================

Panel data that follow many units over many periods allow researchers to
separate three sources of dependence that are confounded in a single cross
section or time series: state dependence, captured by lags of the outcome;
spatial or network interactions, captured by spatially weighted averages of
the outcomes of other units; and common shocks, such as business cycles,
financial crises, or regulatory changes, that hit all units with
heterogeneous intensities. The last source is commonly modelled by a factor
structure $\bm\gamma_i^\top\bm f_t$, also known as interactive effects
\citep{Pesaran:2006,Bai:2003}.

Combining all three features is demanding. Quasi-maximum likelihood
estimators \citep{Yu+DeJong+Lee:2008,Shi+Lee:2017,Bai+Li:2021} are
non-linear, require bias correction for incidental parameters, and do not
accommodate covariates that are endogenous with respect to the idiosyncratic
error. \citet{Cui+Sarafidis+Yamagata:2023} proposed a two-stage
instrumental variables (2SIV) estimator that avoids all three problems: it
is linear in the parameters, free from asymptotic bias, and allows external
instruments. \citet{Chen+Cui+Sarafidis:2025} extended the approach to
heterogeneous slopes with a mean-group IV (MGIV) estimator.
\citet{Kripfganz+Sarafidis:2025} implemented both estimators in the
\proglang{Stata} command \code{spxtivdfreg} and noted that the approach had
not been implemented in any other statistical software.

This article introduces \pkg{spdfiv}, a \proglang{Python} \citep{Python} package
that fills this gap. Its dependencies are limited to \pkg{NumPy}
\citep{Harris+Millman+vanderWalt:2020}, \pkg{SciPy}
\citep{Virtanen+Gommers+Oliphant:2020}, and \pkg{pandas}
\citep{McKinney:2010}, so it is available to every user of the scientific
\proglang{Python} stack at no cost. Spatial econometric software in
\proglang{Python}, notably the \pkg{spreg} package of the \pkg{PySAL}
ecosystem \citep{Rey+Anselin:2010,Rey+Anselin+Amaral:2022}, does not
provide estimators with interactive effects; neither do \pkg{splm}
\citep{Millo+Piras:2012} in \proglang{R} or \pkg{xsmle}
\citep{Belotti+Hughes+Mortari:2017} in \proglang{Stata}.

A translation of an existing command would be of limited interest on its
own. The package therefore also addresses four limitations that
\citet{Kripfganz+Sarafidis:2025} acknowledge in their own discussion.

\begin{enumerate}
\item \emph{Alignment of the spatial weights matrix.} \code{spxtivdfreg}
  does not check that the rows of the weights matrix are ordered like the
  units in the data, leaving this to the user. \pkg{spdfiv} matches the
  matrix to the panel by unit label and refuses unlabelled input unless the
  user explicitly accepts positional matching.
\item \emph{Inference on impacts.} Standard errors of direct, indirect, and
  total impacts are available only through the Delta method. Long-run
  multipliers such as $\beta/(1-\rho-\psi)$ are highly non-linear when the
  process is persistent. \pkg{spdfiv} adds Krinsky--Robb simulation
  \citep{Krinsky+Robb:1986}, the approach of \citet{LeSage+Pace:2009}, and
  computes impacts through a Schur decomposition that remains exact when the
  weights matrix is not diagonalizable.
\item \emph{Outliers in mean-group estimation.} In the application of
  \citet{Kripfganz+Sarafidis:2025}, the MG estimates differ from the pooled
  ones by an order of magnitude for some coefficients, and the authors
  conjecture that a few outlying units are responsible, leaving the question
  for future work. \pkg{spdfiv} provides tools to answer it: influence
  statistics for each unit, outlier flags, and robust summaries of the
  distribution of unit-specific estimates with bootstrap inference.
\item \emph{Unbalanced panels.} \code{spxtivdfreg} requires a balanced
  panel. \pkg{spdfiv} accepts panels with staggered entry or sporadic gaps
  through an explicit, experimental option, documented together with its
  limitations and with simulation evidence of its performance.
\end{enumerate}

The package is validated in three ways. A test suite re-implements the 2SIV
estimator literally from its defining equations and checks that the package
reproduces estimates, variance matrix, and $J$ statistic to machine
precision; Monte Carlo experiments document finite-sample performance; and,
on the bank data of \citet{Kripfganz+Sarafidis:2025}, the package
reproduces their published 2SIV and mean-group estimates, standard errors,
$J$ statistic, and long-run impacts in every printed digit. This
cross-software check uncovered five implementation details of
\code{spxtivdfreg} that are not visible from the equations of the article
but change the results substantially; they are the
defaults of \pkg{spdfiv} and are documented in Section~\ref{sec:application}.

Section~\ref{sec:models} reviews the estimators and describes the
extensions. Section~\ref{sec:package} presents the package,
Section~\ref{sec:illustration} illustrates its use,
Section~\ref{sec:mc} reports Monte Carlo evidence,
Section~\ref{sec:application} applies the package to the bank data,
Section~\ref{sec:timing} reports computing times, and
Section~\ref{sec:summary} concludes.

%% ============================================================================
\section{Models and methods} \label{sec:models}
%% ============================================================================

\subsection{Model}

For units $i=1,\dots,N$ and periods $t=1,\dots,T$, consider
\begin{equation}
y_{it} = \sum_{\tau=0}^{S_\psi}\psi_\tau\sum_{j=1}^N w_{ij}y_{j,t-\tau}
       + \sum_{\tau=1}^{S_\rho}\rho_\tau y_{i,t-\tau}
       + \bm x_{it}^\top\bm\beta + \sum_{j=1}^N w_{ij}\bm x_{jt}^\top\bm\delta
       + \alpha_i + \bm\gamma_{i}^\top\bm f_{t} + \varepsilon_{it},
\label{eq:model}
\end{equation}
where $w_{ij}$ are the elements of a known $N\times N$ spatial weights
matrix $\bm W$ with zero diagonal, $\alpha_i$ are unit effects, $\bm f_t$
are $r_y$ unobserved factors with loadings $\bm\gamma_i$, and
$\varepsilon_{it}$ is an idiosyncratic error. The covariates follow
\begin{equation}
\bm x_{it} = \bm\Gamma_{i}^\top \bm f_{x,t} + \bm v_{it},
\label{eq:xdgp}
\end{equation}
so that they are correlated with the factor component of the error. This
is model~(23) of \citet{Kripfganz+Sarafidis:2025}; with $S_\psi=0$,
$S_\rho=1$ and $\bm\delta=\bm 0$ it reduces to the model of
\citet{Cui+Sarafidis+Yamagata:2023}. Stacking the $T$ observations of unit
$i$ gives $\bm y_i = \bm C_i\bm\theta + \bm u_i$, where $\bm C_i$ collects
the regressors and $\bm u_i = \bm F\bm\gamma_i + \bm\varepsilon_i$.

\subsection{Two-stage IV estimation}

Let $\hat{\bm F}_{x,-\tau}$ be $\sqrt{T}$ times the eigenvectors of the
$r_x$ largest eigenvalues of $(NT)^{-1}\sum_i \bm X_{i,-\tau}\bm X_{i,-\tau}^\top$,
where $\bm X_{i,-\tau}$ holds the instrumenting covariates lagged $\tau$
periods, and let $\bm M_{\hat F_{x,-\tau}}$ be the projection that removes
them. The instrument matrix is
\begin{equation}
\hat{\bm Z}_i = \left(\bm M_{\hat F_{x,-\tau}}\bm X_{i,-\tau},\;
  \textstyle\sum_j w_{ij}\bm M_{\hat F_{x,-\tau}}\bm X_{j,-\tau}\right)_{\tau=0,\dots,L}.
\label{eq:Z}
\end{equation}
The first-stage estimator is the IV estimator
$\hat{\bm\theta} = (\hat{\bm A}^\top\hat{\bm B}^{-1}\hat{\bm A})^{-1}\hat{\bm A}^\top\hat{\bm B}^{-1}\hat{\bm c}$
with $\hat{\bm A}=\sum_i\hat{\bm Z}_i^\top\bm C_i$,
$\hat{\bm B}=\sum_i\hat{\bm Z}_i^\top\hat{\bm Z}_i$, and
$\hat{\bm c}=\sum_i\hat{\bm Z}_i^\top\bm y_i$. In the second stage the
factors $\hat{\bm F}_y$ are estimated from the first-stage residuals
$\hat{\bm u}_i$, projected out by $\bm M_{\hat F_y}$, and
\begin{equation}
\tilde{\bm\theta} = (\tilde{\bm A}^\top\tilde{\bm B}^{-1}\tilde{\bm A})^{-1}\tilde{\bm A}^\top\tilde{\bm B}^{-1}\tilde{\bm c},\quad
\tilde{\bm A} = \sum_i\hat{\bm Z}_i^\top\bm M_{\hat F_y}\bm C_i,\;
\tilde{\bm B} = \sum_i\hat{\bm Z}_i^\top\bm M_{\hat F_y}\hat{\bm u}_i\hat{\bm u}_i^\top\bm M_{\hat F_y}\hat{\bm Z}_i .
\label{eq:2siv}
\end{equation}
\citet{Cui+Sarafidis+Yamagata:2023} show that $\tilde{\bm\theta}$ is
consistent, asymptotically normal, and correctly centred as
$N,T\to\infty$ with $N/T\to c\in(0,\infty)$. The package uses the robust
variance
\begin{equation}
\hat{\bm V} = (\tilde{\bm A}^\top\tilde{\bm B}^{-1}\tilde{\bm A})^{-1}\tilde{\bm A}^\top\tilde{\bm B}^{-1}\hat{\bm\Omega}\tilde{\bm B}^{-1}\tilde{\bm A}(\tilde{\bm A}^\top\tilde{\bm B}^{-1}\tilde{\bm A})^{-1},\qquad
\hat{\bm\Omega}=\sum_i\hat{\bm Z}_i^\top\tilde{\bm e}_i\tilde{\bm e}_i^\top\hat{\bm Z}_i,
\label{eq:V}
\end{equation}
which allows heteroskedasticity and serial correlation within units. By
default $\hat{\bm\Omega}$ is evaluated at the first-stage residuals,
$\hat{\bm\Omega}=\tilde{\bm B}$, as in Equation~(11) of
\citet{Kripfganz+Sarafidis:2025} and in \code{spxtivdfreg}, so that
$\hat{\bm V}=(\tilde{\bm A}^\top\tilde{\bm B}^{-1}\tilde{\bm A})^{-1}$; with
\code{omega="second"} it is evaluated at the second-stage residuals
$\tilde{\bm e}_i=\bm M_{\hat F_y}(\bm y_i-\bm C_i\tilde{\bm\theta})$. Both choices
are consistent. The package reports the overidentification statistic
$J=\bar{\bm g}^\top\hat{\bm\Omega}^{-1}\bar{\bm g}$ with
$\bar{\bm g}=\sum_i\hat{\bm Z}_i^\top\tilde{\bm e}_i$, asymptotically
$\chi^2$ with as many degrees of freedom as there are overidentifying
restrictions. The numbers of factors are chosen by the eigenvalue-ratio
criterion of \citet{Ahn+Horenstein:2013}, which selects between one and
\code{factmax} factors; with \code{zero\_factors=True} their mock eigenvalue
is added so that zero factors can be selected. Unit effects are removed by a
within transformation of all variables, including the instruments, before
the factors are extracted.

\subsection{Mean-group IV estimation}

When $\bm\theta_i=\bm\theta+\bm\eta_i$ with $\bm\eta_i$ independent across
units, \citet{Chen+Cui+Sarafidis:2025} estimate each $\bm\theta_i$ by IV
using the instruments~\eqref{eq:Z}, where the lagged blocks are
additionally projected on the contemporaneous factors (``double
defactoring''), and average:
$\hat{\bm\theta}_{\mathrm{MG}} = N^{-1}\sum_i\hat{\bm\theta}_i$ with variance
estimator $\hat{\bm\Sigma}_\eta/N$, where
$\hat{\bm\Sigma}_\eta=(N-1)^{-1}\sum_i(\hat{\bm\theta}_i-\hat{\bm\theta}_{\mathrm{MG}})(\hat{\bm\theta}_i-\hat{\bm\theta}_{\mathrm{MG}})^\top$
\citep[see also][]{Pesaran+Smith:1995}.

\subsection{Impacts and their inference} \label{sec:impacts}

For covariate $\ell$ the matrix of long-run partial derivatives is
\begin{equation}
\bm S_\ell = \left[a\bm I_N - b\bm W\right]^{-1}(\beta_\ell\bm I_N+\delta_\ell\bm W),\qquad
a = 1-\textstyle\sum_\tau\rho_\tau,\quad b=\sum_\tau\psi_\tau,
\label{eq:S}
\end{equation}
and the short-run matrix sets $a=1$ and $b=\psi_0$. The average direct
impact is $\mathrm{tr}(\bm S_\ell)/N$, the average total impact is
$\bm 1^\top\bm S_\ell\bm 1/N$, and the average indirect impact is their
difference \citep{LeSage+Pace:2009,Debarsy+Ertur+LeSage:2012}.

\paragraph{Computation.} Let $\bm W=\bm Q\bm R\bm Q^*$ be the complex Schur
decomposition, with $\bm Q$ unitary and $\bm R$ upper triangular with the
eigenvalues $\lambda_1,\dots,\lambda_N$ on its diagonal. Because $\bm S_\ell$
is a rational function of $\bm W$, its trace equals
$\sum_j(\beta_\ell+\delta_\ell\lambda_j)/(a-b\lambda_j)$ \emph{whether or not
$\bm W$ is diagonalizable}, and
$\bm 1^\top\bm S_\ell\bm 1 = (\bm 1^\top\bm Q)(a\bm I-b\bm R)^{-1}\bm Q^*(\beta_\ell\bm 1+\delta_\ell\bm W\bm 1)$
requires a single triangular solve. When all row sums of $\bm W$ equal $s$,
as after row standardization, the total impact reduces to
$(\beta_\ell+\delta_\ell s)/(a-bs)$. After one $O(N^3)$ decomposition, each
evaluation costs $O(N)$ for direct impacts and at most $O(N^2)$ for total
impacts, which makes simulation-based inference cheap. The Schur form
matters in practice: asymmetric $k$-nearest-neighbour matrices, common in
applied work, are often numerically close to defective, and in our
experiments the eigenvector matrices of such matrices had condition numbers
above $10^{11}$, which makes formulas based on $\bm W=\bm V\bm\Lambda\bm V^{-1}$
unreliable.

\paragraph{Inference.} Impacts are smooth functions $h(\bm\theta)$ of the
coefficients. Besides the Delta method, \pkg{spdfiv} offers
Krinsky--Robb simulation: draws $\bm\theta^{(r)}\sim
N(\tilde{\bm\theta},\hat{\bm V})$ are mapped into $h(\bm\theta^{(r)})$, and
percentile intervals are reported. Draws that violate the stability
conditions $\psi_0\omega<1$ and $a-b\omega>0$, with $\omega$ the spectral
radius of $\bm W$, are discarded and their share is reported. The two
methods share the same first-order justification; the simulation intervals
are not forced to be symmetric, which matters when $a-b\omega$ is close to
zero and $h$ is strongly convex.

\subsection{Diagnostics for heterogeneous slopes} \label{sec:diag}

Under the random-coefficient assumption the estimates $\hat{\bm\theta}_i$
are, up to estimation error, independent draws from the distribution of
$\bm\theta_i$, and the MG estimator is the sample mean of that
distribution. A sample mean is sensitive to a few extreme draws, which may
reflect data errors, structural breaks, or units with nearly
uninformative instruments. \pkg{spdfiv} offers three tools.

\begin{enumerate}
\item \emph{Influence.} Deleting unit $i$ changes the MG estimate by
  $\hat{\bm\theta}_{\mathrm{MG}}-\hat{\bm\theta}_{\mathrm{MG}(-i)}=(\hat{\bm\theta}_i-\hat{\bm\theta}_{\mathrm{MG}})/(N-1)$.
  The package reports this change in units of the MG standard error.
\item \emph{Outlier flags.} The modified $z$-score
  $0.6745(\hat\theta_{i\ell}-\mathrm{med}_\ell)/\mathrm{MAD}_\ell$ of
  \citet{Iglewicz+Hoaglin:1993}, with the median and the median absolute
  deviation computed across units, flags estimates with $|z|>3.5$.
\item \emph{Robust summaries.} Coefficient-wise trimmed means and medians
  \citep{Huber+Ronchetti:2009}, with standard errors and percentile
  intervals from a nonparametric bootstrap that resamples units
  \citep{Efron+Tibshirani:1993}; resampling units is justified by the
  independence of $\bm\eta_i$ across units.
\end{enumerate}

The influence table also reports, for each unit, the condition number of
$\hat{\bm Z}_i^\top\hat{\bm Z}_i$; very large values mark units whose
instruments are (nearly) collinear, for example because a covariate does not
vary over time for that unit, so that their estimates depend on numerical
tolerances rather than on the data.

Two caveats are stated in the documentation and bear repeating. The
trimmed mean and the median estimate the centre of the distribution of
$\bm\theta_i$, which equals the mean $\bm\theta$ only if the distribution is
symmetric. And the tools are diagnostics: they show which units drive an
estimate, but whether those units should be excluded is a substantive
decision for the researcher.

\subsection{Unbalanced panels (experimental)} \label{sec:unbal}

Two departures from balance are supported. With \emph{staggered} entry or
exit, each unit is observed over one contiguous spell; with \emph{sporadic}
gaps, interior observations may also be missing. In both cases three
ingredients change.

\begin{enumerate}
\item The spatial lag in period $t$ uses the neighbours available in that
  period. By default the remaining weights in each row are rescaled to keep
  the original row sum, so that a balanced panel is analysed with exactly
  $\bm W$; the alternative, dropping absent neighbours without rescaling, is
  available because, as \citet{Kelejian+Prucha:2010} stress, normalization
  is a modelling choice.
\item Factors are extracted from the $T\times T$ matrix whose $(s,t)$
  element averages over the units observed in both periods, the
  ``all-purpose'' estimator of \citet{Xiong+Pelger:2023}; for a balanced
  panel it coincides with the usual matrix.
\item The projection $\bm M_{\hat F}$ for unit $i$ uses only the rows of
  $\hat{\bm F}$ for the periods in which unit $i$ is observed.
\end{enumerate}

The theory of \citet{Cui+Sarafidis+Yamagata:2023} and
\citet{Chen+Cui+Sarafidis:2025} assumes a balanced panel, and extending it
is beyond the scope of a software article. The option is therefore labelled
experimental, triggers a warning, and is evaluated by simulation in
Section~\ref{sec:mc}. The two cases differ in substance. Staggered entry is
consistent with a model in which a unit does not exist before it enters, so
that the period-specific weights are the correct ones. Sporadic gaps in
units that do exist are a missing-data problem: the spatial lag computed
from observed neighbours then differs from the true one, which can bias the
spatial coefficient \citep[cf.][]{Wang+Lee:2013}.

%% ============================================================================
\section[The spdfiv package]{The \pkg{spdfiv} package} \label{sec:package}
%% ============================================================================

\subsection{Design}

The package is written in pure \proglang{Python}. All computations are
vectorized with \pkg{NumPy} or reduce to per-unit linear algebra on small
matrices, so the cost grows linearly in $N$ apart from the eigen
decompositions of $T\times T$ matrices for the factors and one Schur
decomposition of $\bm W$ for impacts. The main entry point is
\fct{fit}; it returns a results object with methods for impacts and
heterogeneity diagnostics. Table~\ref{tab:options} maps its arguments to the
options of \code{spxtivdfreg}, to ease migration and cross-checks.

\begin{table}[t!]
\centering\small
\begin{tabular}{lll}
\toprule
\pkg{spdfiv} argument & \code{spxtivdfreg} option & Meaning \\
\midrule
\code{W=} (labelled) & \code{spmatrix()} & spatial weights matrix \\
\code{splag=True} & \code{splag} & spatial lag of $y$ \\
\code{tlags=}$p$ & \code{tlags(}$p$\code{)} & time lags of $y$ \\
\code{sptlags=}$q$ & \code{sptlags(}$q$\code{)} & spatial time lags of $y$ \\
\code{spx=[...]} & \code{spindepvars()} & spatial lags of covariates \\
\code{iv=[...]} & \code{iv(}\textit{varlist}\code{)} & exogenous instrumenting variables \\
\code{iv\_lags=}$L$ & \code{iv(..., lags(}$L$\code{))} & lags of instruments \\
\code{iv\_splags=True} & \code{iv(..., splags)} & spatial lags of instruments \\
\code{factmax=}$k$ & \code{factmax(}$k$\code{)} & maximum number of factors \\
\code{eigratio=False} & \code{noeigratio} & use exactly \code{factmax} factors \\
\code{zero\_factors=True} & --- & allow zero factors (mock eigenvalue) \\
\code{common\_rank=True} & (always) & same number of factors for all lag blocks \\
\code{omega="first"} & (always) & $\hat{\bm\Omega}$ from first-stage residuals \\
\code{std=True} & \code{std} & standardize within periods before factor extraction \\
\code{absorb=True} & \code{absorb(}\textit{id}\code{)} & unit fixed effects \\
\code{estimator="fstage"} & \code{fstage} & first-stage IV \\
\code{estimator="mg"} & \code{mg} & mean-group IV \\
\code{doubledefact=} & \code{doubledefact} & double defactoring of lagged instruments \\
\code{unbalanced=} & --- & staggered or sporadic missingness (experimental) \\
\code{renormalize=} & --- & row rescaling with absent neighbours \\
\bottomrule
\end{tabular}
\caption{Arguments of \fct{fit} and corresponding options of the
\proglang{Stata} command \code{spxtivdfreg}.}
\label{tab:options}
\end{table}

\subsection{Results}

{\sloppy The object returned by \fct{fit} stores the coefficients (\code{params}),
their variance matrix (\code{cov}), standard errors (\code{bse}), $z$
statistics (\code{zvalues}), $p$ values (\code{pvalues}), confidence
intervals (\code{conf\_int()}), the numbers of factors
(\code{factors\_x}, one per instrument lag block, and \code{factors\_u}),
the $J$ statistic with its degrees of freedom and $p$ value, and the
decomposition of the error variance into factor and idiosyncratic parts.
For the MG estimator, \code{unit\_params} and \code{unit\_se} hold the
unit-specific estimates. The methods \fct{impacts}, \fct{mg\_influence},
and \fct{mg\_robust} implement Sections~\ref{sec:impacts}
and~\ref{sec:diag}, and \fct{summary} prints a coefficient table.\par}

\subsection{Alignment of the weights matrix}

The weights matrix is passed as a \pkg{pandas} \code{DataFrame} whose
index and columns are unit identifiers, or as an array together with the
identifiers (\code{W\_ids}). The package checks that the labels on both
axes agree, that they coincide with the units in the data, that there are
no missing values, and that the diagonal is zero; it then reorders the
matrix to the internal ordering of units. An unlabelled array is rejected
unless \code{assume\_ordered=True}. The helper \fct{read\_weights} reads a
labelled matrix from a CSV or Excel file.

\subsection{Validation}

The test suite, run with \code{python -m unittest}, contains a literal
implementation of equations (5)--(12) of \citet{Kripfganz+Sarafidis:2025},
for both choices of $\hat{\bm\Omega}$, written with explicit $T\times T$
projection matrices and independent of the package code. The package reproduces its first- and second-stage estimates
to an absolute tolerance of $10^{-10}$, its variance matrix to a relative
tolerance of $10^{-8}$, and its $J$ statistic to eight decimal places. With
no factors and no lags the first-stage estimator coincides with two-stage
least squares, which is also tested. Further tests check that permuting a
labelled weights matrix leaves the estimates unchanged, that misaligned or
invalid matrices are rejected, that the Schur-based impacts agree with dense
matrix inversion for both row-standardized and spectrally normalized
matrices, that the balanced-panel estimates are unchanged by the unbalanced
option, and that the influence statistic equals the leave-one-out change in
the MG estimate.

%% ============================================================================
\section{Illustration} \label{sec:illustration}
%% ============================================================================

This section uses simulated data so that the true parameters are known and
the output can be reproduced exactly. The data-generating process is
model~\eqref{eq:model} with $\psi_0=0.3$, $\rho_1=0.4$,
$\bm\beta=(1,0.5)^\top$, two AR(1) factors that enter both the error and the
covariates, loadings of the covariates correlated with those of the error,
heteroskedastic errors, unit effects, and a row-standardized six-nearest-neighbour
weights matrix; it is available as \fct{simulate\_panel}.

\begin{CodeInput}
>>> import pandas as pd
>>> import spdfiv as sd
>>> data, W, truth = sd.simulate_panel(N=300, T=40, psi=0.3, rho=0.4,
...                                    beta=(1.0, 0.5), seed=2025)
>>> print(data.head(3)); print(W.shape, W.index[:5].tolist())
\end{CodeInput}
\VerbatimInput[fontsize=\small]{../replication/results/illustration_1.txt}

The data are in long format; \code{W} is a labelled \code{DataFrame}. The
2SIV estimator with a spatial lag, one time lag, and instruments formed from
the defactored covariates, their first lags, and the spatial lags of both,
is obtained with

\begin{CodeInput}
>>> res = sd.fit(data, y="y", x=["x1", "x2"], unit="id", time="time", W=W,
...              splag=True, tlags=1, iv_lags=1, iv_splags=True)
>>> print(res.summary())
\end{CodeInput}
\VerbatimInput[fontsize=\small]{../replication/results/illustration_2.txt}

The eigenvalue-ratio criterion selects the true number of factors in each
instrument block and in the residuals, the estimates are close to the true
values, and the $J$ test does not reject. Ignoring the factors
(\code{factmax=0}) yields IV estimates based on raw covariates:

\begin{CodeInput}
>>> res0 = sd.fit(data, "y", ["x1", "x2"], "id", "time", W, tlags=1,
...               iv_lags=1, factmax=0)
\end{CodeInput}
\VerbatimInput[fontsize=\small]{../replication/results/illustration_3.txt}

The slope estimates are biased upwards by 0.34 to 0.38 because the
covariates are correlated with the omitted factor component, and the $J$
test rejects at the 5\% level.

\paragraph{Impacts.} Long-run impacts with simulation-based inference are
requested with

\begin{CodeInput}
>>> lr_delta = res.impacts("lr", method="delta")
>>> lr_sim = res.impacts("lr", method="simulation", draws=5000, seed=1)
>>> print(lr_sim.round(4))
>>> print(pd.concat({"delta": lr_delta[["se", "lower", "upper"]],
...                  "simulation": lr_sim[["se", "lower", "upper"]]},
...                 axis=1).loc["total"].round(4))
\end{CodeInput}
\VerbatimInput[fontsize=\small]{../replication/results/illustration_4.txt}

The second panel compares the two methods for the total impacts. The
simulation intervals are shifted to the right of the symmetric Delta
intervals, reflecting the convexity of $\beta/(1-\rho-\psi)$ in the
persistence parameters. The true long-run total impacts are $1/0.3=3.333$
for \code{x1} and $1.667$ for \code{x2}; both intervals cover them.

\paragraph{Safeguards.} Passing the weights as an unlabelled array is
refused, and reversing the order of a labelled matrix leaves the estimates
unchanged:

\begin{CodeInput}
>>> sd.fit(data, "y", ["x1", "x2"], "id", "time", W.to_numpy())
>>> perm = W.index[::-1]
>>> res_p = sd.fit(data, "y", ["x1", "x2"], "id", "time", W.loc[perm, perm],
...                tlags=1, iv_lags=1)
\end{CodeInput}
\VerbatimInput[fontsize=\small]{../replication/results/illustration_5.txt}

\paragraph{Heterogeneous slopes.} The next data set has heterogeneous
slopes, centred on the same values, and 3\% of units (nine units) whose
covariate slopes are shifted by $+3$, mimicking units with data problems or
a different structure.

\begin{CodeInput}
>>> dh, Wh, th = sd.simulate_panel(N=300, T=40, het=0.2, outlier_share=0.03,
...                                seed=7)
>>> mg = sd.fit(dh, "y", ["x1", "x2"], "id", "time", Wh, tlags=1,
...             iv_lags=1, estimator="mg")
>>> print(mg.summary())
\end{CodeInput}
\VerbatimInput[fontsize=\small]{../replication/results/illustration_6.txt}

The MG estimates of the covariate slopes, 1.072 and 0.598, lie well above
the centre of the regular units, 1 and 0.5. The influence diagnostics show
why:

\begin{CodeInput}
>>> infl = sd.mg_influence(mg, threshold=3.5)
>>> cols = ["unit", "parameter", "theta_i", "mod_z", "dfmg", "flag"]
>>> print(infl.sort_values("dfmg", key=abs, ascending=False)[cols].head(6))
>>> flagged = sd.flagged_units(infl, parameters=["x1", "x2"])
\end{CodeInput}
\VerbatimInput[fontsize=\small]{../replication/results/illustration_7.txt}

Each of the most influential units moves an MG coefficient by about a third
of its standard error. Flagging on the covariate slopes recovers all nine
aberrant units and one regular unit. The two units flagged for the spatial
and time lag coefficients have imprecise unit-level estimates; with $T=40$
per-unit IV estimates of these coefficients are noisy, which is why
flagging is best restricted to the parameters of interest. Robust summaries
and the MG estimate without the flagged units are

\begin{CodeInput}
>>> print(sd.mg_robust(mg, method="trimmed", trim=0.1, bootstrap=999, seed=1))
>>> print(sd.mg_robust(mg, method="mean", exclude=flagged, bootstrap=999,
...                    seed=1))
\end{CodeInput}
\VerbatimInput[fontsize=\small]{../replication/results/illustration_8.txt}

Both the trimmed mean and the MG estimate without the flagged units are close
to the centre of the regular units, with standard errors about half those of
the full MG estimate.

\paragraph{Staggered entry.} Finally, 30\% of the units enter the panel
at random dates during the first half of the sample.

\begin{CodeInput}
>>> du, Wu, tu = sd.simulate_panel(N=300, T=40, entry_share=0.3, seed=11)
>>> ru = sd.fit(du, "y", ["x1", "x2"], "id", "time", Wu, tlags=1,
...             iv_lags=1, unbalanced="staggered")
\end{CodeInput}
\VerbatimInput[fontsize=\small]{../replication/results/illustration_9.txt}

The first lines show how many units start in each of the first periods. The
estimator uses all 10{,}800 available observations and issues a warning that
the option is experimental.

%% ============================================================================
\section{Monte Carlo evidence} \label{sec:mc}
%% ============================================================================

All experiments use the data-generating process of
Section~\ref{sec:illustration}, 1{,}000 replications, a new weights matrix
and new loadings in each replication, and the specification with one time
lag, one instrument lag, and spatial lags of the instruments (eight
instruments, four parameters). The code is in
\code{replication/monte\_carlo.py}.

\subsection{Homogeneous slopes}

Table~\ref{tab:mc1} compares the 2SIV estimator with the first-stage IV
estimator and with IV estimation that ignores the factors
(\code{factmax=0}), for $(N,T)\in\{(50,50),(100,100),(200,200)\}$. Size
refers to two-sided $t$ tests at the 5\% level, and the last column reports
the rejection frequency of the $J$ test at the 5\% level.

\begin{table}[t!]
\centering\scriptsize
\begin{tabular}{llrrrrrrrrrr}
\toprule
 & & \multicolumn{3}{c}{$\psi_0$} & \multicolumn{3}{c}{$\rho_1$} & \multicolumn{3}{c}{$\beta_1$} & \\
\cmidrule(lr){3-5}\cmidrule(lr){6-8}\cmidrule(lr){9-11}
$(N,T)$ & Estimator & Bias & RMSE & Size & Bias & RMSE & Size & Bias & RMSE & Size & $J$ \\
\midrule
$(50,50)$ & 2SIV & 0.005 & 0.034 & 9.1 & $-$0.003 & 0.020 & 11.1 & 0.002 & 0.026 & 10.4 & 4.9 \\
 & First-stage IV & 0.021 & 0.050 & 12.3 & $-$0.018 & 0.030 & 8.1 & 0.012 & 0.029 & 7.5 & 18.0 \\
 & IV, no factors & 0.028 & 0.036 & 41.6 & $-$0.010 & 0.016 & 30.4 & 0.354 & 0.357 & 100.0 & 21.5 \\
\addlinespace
$(100,100)$ & 2SIV & 0.001 & 0.015 & 6.5 & $-$0.001 & 0.009 & 9.6 & 0.000 & 0.012 & 7.6 & 3.5 \\
 & First-stage IV & 0.003 & 0.018 & 6.6 & $-$0.009 & 0.016 & 7.5 & 0.005 & 0.014 & 6.5 & 14.8 \\
 & IV, no factors & 0.019 & 0.023 & 45.6 & $-$0.007 & 0.010 & 40.6 & 0.360 & 0.361 & 100.0 & 54.0 \\
\addlinespace
$(200,200)$ & 2SIV & 0.000 & 0.007 & 5.8 & 0.000 & 0.004 & 5.3 & 0.000 & 0.006 & 4.5 & 3.8 \\
 & First-stage IV & 0.001 & 0.008 & 6.6 & $-$0.004 & 0.007 & 5.5 & 0.003 & 0.007 & 6.2 & 17.4 \\
 & IV, no factors & 0.012 & 0.015 & 49.2 & $-$0.004 & 0.007 & 46.1 & 0.363 & 0.363 & 100.0 & 91.1 \\
\bottomrule
\end{tabular}

\caption{Homogeneous slopes: bias, root mean squared error, and size (\%) of
5\% $t$ tests for $\psi_0=0.3$, $\rho_1=0.4$, $\beta_1=1$, and rejection
frequency (\%) of the $J$ test. 1{,}000 replications.}
\label{tab:mc1}
\end{table}

The 2SIV estimator is practically unbiased in all designs, and its RMSE
halves as $N$ and $T$ double. Its $t$ tests over-reject somewhat in the
smallest design (sizes between 9.1\% and 11.1\%), but their size is between
4.5\% and 5.8\% when $N=T=200$, and the $J$ test has size between 3.5\% and
4.9\%. The eigenvalue-ratio criterion selects the true number of factors in
the instruments in at least 99.8\% of the replications and in the residuals
in at least 97.2\%. The first-stage estimator is consistent but less
precise; its $J$ test rejects in 15\%--18\% of the replications because it
ignores the factors in the error, which illustrates why the second stage is
the default. Ignoring the factors produces a bias of about 0.36 in
$\beta_1$ that does not vanish as the sample grows, so that $t$ tests on
$\beta_1$ always reject. Reassuringly, the $J$ test detects this
misspecification with a frequency that rises from 21.5\% to 91.1\%.

\subsection{Inference on long-run impacts}

To make the long-run multipliers strongly non-linear, Table~\ref{tab:mc2}
uses $\psi_0=0.35$ and $\rho_1=0.5$, so that $1-\rho_1-\psi_0=0.15$ and the
total long-run impact of $x_1$ is $6.67$. It reports the coverage of
nominal 95\% intervals, the frequencies with which the true value lies below
or above the interval, and the median interval width. In the smallest
design, the point estimates violated the stability condition in 22 of the
1{,}000 replications, for which impacts are undefined; these replications are
excluded.

\begin{table}[t!]
\centering\scriptsize
\begin{tabular}{llrrrrrrrr}
\toprule
 & & \multicolumn{4}{c}{Delta method} & \multicolumn{4}{c}{Simulation (Krinsky--Robb)} \\
\cmidrule(lr){3-6}\cmidrule(lr){7-10}
$(N,T)$ & Impact & Cover & Below & Above & Width & Cover & Below & Above & Width \\
\midrule
$(30,30)$ & direct & 88.7 & 1.7 & 9.6 & 0.63 & 87.2 & 7.1 & 5.7 & 0.71 \\
 & indirect & 89.1 & 0.2 & 10.7 & 5.60 & 86.9 & 7.6 & 5.5 & 7.15 \\
 & total & 88.7 & 0.2 & 11.1 & 6.05 & 87.1 & 7.5 & 5.4 & 7.60 \\
\addlinespace
$(50,50)$ & direct & 92.7 & 1.6 & 5.7 & 0.40 & 90.4 & 5.6 & 4.0 & 0.41 \\
 & indirect & 90.8 & 0.5 & 8.7 & 3.76 & 88.6 & 6.4 & 5.0 & 4.15 \\
 & total & 90.8 & 0.5 & 8.7 & 4.08 & 88.3 & 6.4 & 5.3 & 4.48 \\
\addlinespace
$(100,100)$ & direct & 92.9 & 2.2 & 4.9 & 0.19 & 93.7 & 3.2 & 3.1 & 0.19 \\
 & indirect & 93.3 & 1.4 & 5.3 & 1.85 & 94.5 & 2.7 & 2.8 & 1.89 \\
 & total & 93.5 & 1.4 & 5.1 & 2.00 & 94.2 & 2.6 & 3.2 & 2.04 \\
\bottomrule
\end{tabular}

\caption{Long-run impacts of $x_1$ with $\psi_0=0.35$, $\rho_1=0.5$:
coverage (\%) of nominal 95\% intervals, frequency (\%) with which the true
value lies below the lower or above the upper limit, and median width.}
\label{tab:mc2}
\end{table}

The results give a nuanced picture that users should know about. The Delta
intervals are symmetric while the sampling distribution of the impacts is
skewed to the right; as a result almost all of their non-coverage is on one
side: the true value lies above the upper limit in 9\%--11\% of the
replications for indirect and total impacts when $N,T\le 50$, and below the
lower limit in less than 1\%. Simulation intervals spread the errors almost
evenly between the two tails. Their overall coverage, however, is not
better in small samples (86.9\%--88.6\% compared with 88.7\%--90.8\% for the Delta
method), because both methods inherit the same first-order approximation of
the variance. In the largest design both methods are close to nominal
coverage, simulation slightly more so (94.2\% compared with 93.5\% for total
impacts). The practical recommendation is therefore to report simulation
intervals when the process is persistent, because they convey the
asymmetry of the uncertainty, while recognizing that neither method is
exact in small samples.

\subsection{Mean-group estimation with aberrant units}

The data for Table~\ref{tab:mc3} have heterogeneous slopes, with
$\psi_{0i}$ and $\rho_{1i}$ uniform on $\pm0.1$ around their means and
$\beta_{\ell i}$ uniform on $\pm0.4$ around $\beta_\ell$. A share of 0\%,
5\%, or 10\% of the units is \emph{aberrant}: both of their covariate slopes
are shifted by $+3$. The target is the mean of the regular units,
$\beta_1=1$. The table reports the MG estimator, the coefficient-wise 10\%
trimmed mean and the median of the unit-specific estimates (with 499
bootstrap replications over units), and the MG estimator after removing the
units flagged on $\beta_1$ or $\beta_2$, together with the true and false
positive rates of the flags.

\begin{table}[t!]
\centering\scriptsize
{\setlength{\tabcolsep}{4pt}\begin{tabular}{llrrrrrrrrr}
\toprule
 & & \multicolumn{3}{c}{No aberrant units} & \multicolumn{3}{c}{5\% aberrant} & \multicolumn{3}{c}{10\% aberrant} \\
\cmidrule(lr){3-5}\cmidrule(lr){6-8}\cmidrule(lr){9-11}
$(N,T)$ & Estimator & Bias & RMSE & Cover & Bias & RMSE & Cover & Bias & RMSE & Cover \\
\midrule
$(100,50)$ & MG & 0.001 & 0.031 & 94.5 & 0.152 & 0.155 & 36.3 & 0.300 & 0.302 & 0.0 \\
 & Trimmed MG (10\%) & 0.002 & 0.031 & 95.1 & 0.031 & 0.044 & 87.8 & 0.065 & 0.073 & 63.7 \\
 & Median & 0.001 & 0.041 & 95.3 & 0.026 & 0.050 & 91.7 & 0.049 & 0.066 & 82.0 \\
 & MG excl.\ flagged & 0.002 & 0.031 & 94.3 & 0.003 & 0.032 & 95.2 & 0.000 & 0.034 & 93.1 \\
 & \emph{Flagging: TPR / FPR} & \multicolumn{3}{c}{-- / 0.3} & \multicolumn{3}{c}{100.0 / 0.2} & \multicolumn{3}{c}{100.0 / 0.1} \\
\addlinespace
$(200,100)$ & MG & 0.000 & 0.019 & 94.2 & 0.150 & 0.151 & 0.0 & 0.300 & 0.300 & 0.0 \\
 & Trimmed MG (10\%) & 0.000 & 0.021 & 94.3 & 0.023 & 0.031 & 81.6 & 0.054 & 0.058 & 40.4 \\
 & Median & $-$0.001 & 0.029 & 94.8 & 0.020 & 0.035 & 89.9 & 0.044 & 0.054 & 70.4 \\
 & MG excl.\ flagged & 0.000 & 0.019 & 94.2 & 0.000 & 0.019 & 94.0 & 0.000 & 0.020 & 94.4 \\
 & \emph{Flagging: TPR / FPR} & \multicolumn{3}{c}{-- / 0.0} & \multicolumn{3}{c}{100.0 / 0.0} & \multicolumn{3}{c}{100.0 / 0.0} \\
\bottomrule
\end{tabular}
}
\caption{Heterogeneous slopes with aberrant units (slopes shifted by $+3$):
bias, RMSE, and coverage (\%) of nominal 95\% intervals for the centre of
the regular units, $\beta_1=1$, and true / false positive rates (\%) of the
outlier flags. 1{,}000 replications.}
\label{tab:mc3}
\end{table}

Without aberrant units all four estimators are centred on the truth and
have coverage between 94.2\% and 95.3\%; the median pays an efficiency price
(RMSE 0.041 against 0.031 for MG when $N=100$), and the flags mark at most
0.3\% of the regular units. With 5\% aberrant units the MG estimator is
biased by 0.15, and its coverage falls to 36\% for $N=100$ and to zero for
$N=200$: a larger sample only makes the wrong answer more precise. The
trimmed mean and the median remove most of the bias but not all of it,
because the symmetric trimming also cuts regular units from the lower tail,
so their coverage remains below nominal (82\%--92\%) and deteriorates when
10\% of the units are aberrant. Excluding the flagged units restores both
the centre and the coverage (93\%--95\%) in all designs.

With a shift of $+3$ the aberrant units are far outside the distribution of
the regular slopes, which makes detection easy: every aberrant unit is
flagged in every design. To see how the tools behave when detection is hard, we repeated the
$(N,T)=(200,100)$ design with a shift of $+1$, only 2.5 times the
half-width of the distribution of regular slopes (1{,}000 replications;
\code{monte\_carlo.py}, experiment \code{3b}). The flags then identify
55.5\% of the aberrant units when 5\% are aberrant and 39.6\% when 10\% are,
with no false positives. Excluding the flagged units still halves the bias
of the MG estimator with 5\% aberrant units (0.020 against 0.050) but no
longer restores coverage (84.6\%), and with 10\% aberrant units it is no
better than the trimmed mean (bias 0.056 and 0.053). The median has the
smallest bias in both cases (0.022 and 0.044). The diagnostics are
therefore most useful when aberrant units are clearly separated, which is
the case they are designed to reveal; when contamination blends into the
regular distribution, robust summaries and exclusion give similar,
partially corrected answers, and the gap between them and the MG estimate
is itself the informative signal.


\subsection{Unbalanced panels}

Table~\ref{tab:mc4} considers $N=T=100$ with two kinds of missing data.
With \emph{staggered entry}, 30\% of the units enter at a random date
during the first half of the sample and, consistently with
Section~\ref{sec:unbal}, do not interact with other units before they
enter. The experimental unbalanced estimator is compared with the two
balanced alternatives a user of existing software would have: dropping the
late units (and re-standardizing the weights matrix of the remaining units)
or keeping all units but only the periods in which all of them are
observed. With \emph{sporadic} gaps, 2\% of the unit-period cells are
missing at random; the only balanced alternative is to keep the units
without gaps.

\begin{table}[t!]
\centering\scriptsize
{\setlength{\tabcolsep}{4pt}\begin{tabular}{llrrrrrrrrrr}
\toprule
 & & & \multicolumn{3}{c}{$\psi_0$} & \multicolumn{3}{c}{$\rho_1$} & \multicolumn{3}{c}{$\beta_1$} \\
\cmidrule(lr){4-6}\cmidrule(lr){7-9}\cmidrule(lr){10-12}
Design & Estimator & $NT$ & Bias & RMSE & Size & Bias & RMSE & Size & Bias & RMSE & Size \\
\midrule
Staggered & Unbalanced, all units & 9149 & 0.001 & 0.015 & 6.6 & 0.001 & 0.009 & 8.1 & 0.003 & 0.013 & 8.6 \\
 & Drop late units & 6930 & $-$0.070 & 0.073 & 97.6 & 0.003 & 0.012 & 10.0 & 0.004 & 0.016 & 9.5 \\
 & Common period & 5112 & 0.002 & 0.023 & 7.9 & $-$0.002 & 0.013 & 8.9 & 0.002 & 0.018 & 8.2 \\
\addlinespace
Sporadic & Unbalanced, all units & 9507 & $-$0.003 & 0.016 & 7.4 & 0.000 & 0.009 & 9.3 & 0.001 & 0.012 & 7.5 \\
 & Complete units only & 1307 & $-$0.222 & 0.231 & 97.1 & 0.009 & 0.037 & 25.2 & 0.019 & 0.055 & 26.8 \\
\bottomrule
\end{tabular}
}
\caption{Unbalanced panels, $N=T=100$: average number of observations used,
bias, RMSE, and size (\%) of 5\% $t$ tests. 1{,}000 replications.}
\label{tab:mc4}
\end{table}

Under staggered entry the unbalanced estimator is centred on the true
values, and its sizes (6.6\%--8.6\%) are close to those of the 2SIV
estimator in balanced panels of the same dimensions in
Table~\ref{tab:mc1}. Dropping the late units is not a harmless loss of
information: it removes nodes from the network and biases the spatial
coefficient by $-0.07$, so that $t$ tests on $\psi_0$ reject in 98\% of the
replications. Restricting the sample to the common period is valid but
discards almost half of the observations, which raises the RMSE of $\psi_0$
from 0.015 to 0.023. With 2\% sporadic gaps the unbalanced estimator is again
nearly unbiased (bias of $\psi_0$ equal to $-0.003$), whereas keeping only
the units without gaps leaves about 14\% of the observations and severely
biases the spatial coefficient. These results support the usefulness of the
option in designs like these. They do not replace a formal theory, and the
sporadic case in particular could behave worse when missingness is more
frequent or related to the outcome; the option therefore remains labelled
experimental.



%% ============================================================================
\section{Empirical application: credit risk in US banks} \label{sec:application}
%% ============================================================================

%% Tables and numbers in this section are generated by
%%   python replication/bank_replication.py v113i06.dta W.csv
%% from the replication materials of Kripfganz and Sarafidis (2025).
\providecommand{\BankMaxRelCoef}{}\renewcommand{\BankMaxRelCoef}{0.000}
\providecommand{\BankMaxRelSE}{}\renewcommand{\BankMaxRelSE}{0.00}
\providecommand{\BankMaxMGdiffSE}{}\renewcommand{\BankMaxMGdiffSE}{0.00}
\providecommand{\BankMedMGdiffSE}{}\renewcommand{\BankMedMGdiffSE}{0.00}
\providecommand{\BankJ}{}\renewcommand{\BankJ}{18.825}
\providecommand{\BankpJ}{}\renewcommand{\BankpJ}{0.468}
\providecommand{\BankNflag}{}\renewcommand{\BankNflag}{95}
\providecommand{\BankNunits}{}\renewcommand{\BankNunits}{350}
\providecommand{\BankNflagSize}{}\renewcommand{\BankNflagSize}{24}
\providecommand{\BankNflagBuffer}{}\renewcommand{\BankNflagBuffer}{15}
\providecommand{\BankNill}{}\renewcommand{\BankNill}{5}
\providecommand{\BankNillFlag}{}\renewcommand{\BankNillFlag}{1}
\providecommand{\BankMGpsi}{}\renewcommand{\BankMGpsi}{0.032}
\providecommand{\BankMGpsiNoDD}{}\renewcommand{\BankMGpsiNoDD}{0.151}
\providecommand{\BankMGsize}{}\renewcommand{\BankMGsize}{2.004}
\providecommand{\BankMGbuffer}{}\renewcommand{\BankMGbuffer}{$-$0.376}
\providecommand{\BankTwoSize}{}\renewcommand{\BankTwoSize}{0.223}
\providecommand{\BankTwoBuffer}{}\renewcommand{\BankTwoBuffer}{$-$0.055}
\providecommand{\BankTwoPsi}{}\renewcommand{\BankTwoPsi}{0.394}
\providecommand{\BankMedSize}{}\renewcommand{\BankMedSize}{0.807}
\providecommand{\BankMedBuffer}{}\renewcommand{\BankMedBuffer}{$-$0.220}
\providecommand{\BankMedPsi}{}\renewcommand{\BankMedPsi}{0.034}
\providecommand{\BankTrimSize}{}\renewcommand{\BankTrimSize}{1.282}
\providecommand{\BankExclSize}{}\renewcommand{\BankExclSize}{0.707}
\providecommand{\BankExclBuffer}{}\renewcommand{\BankExclBuffer}{$-$0.234}
\providecommand{\BankExclPsi}{}\renewcommand{\BankExclPsi}{0.023}
\providecommand{\BankExclIllSize}{}\renewcommand{\BankExclIllSize}{1.997}
\providecommand{\BankGapCutSize}{}\renewcommand{\BankGapCutSize}{67}
\providecommand{\BankGapCutBuffer}{}\renewcommand{\BankGapCutBuffer}{49}
\providecommand{\BankNlate}{}\renewcommand{\BankNlate}{105}
\providecommand{\BankNobsStag}{}\renewcommand{\BankNobsStag}{11{,}602}
\providecommand{\BankNobsDrop}{}\renewcommand{\BankNobsDrop}{8{,}575}
\providecommand{\BankPsiStag}{}\renewcommand{\BankPsiStag}{0.352}
\providecommand{\BankPsiDrop}{}\renewcommand{\BankPsiDrop}{0.367}
\providecommand{\BankMaxStagDiffSE}{}\renewcommand{\BankMaxStagDiffSE}{0.58}
\providecommand{\BankMaxDropDiffSE}{}\renewcommand{\BankMaxDropDiffSE}{1.14}
\providecommand{\BankImpRel}{}\renewcommand{\BankImpRel}{0.00}
\providecommand{\BankImpRelSE}{}\renewcommand{\BankImpRelSE}{0.00}
\providecommand{\BankSimStable}{}\renewcommand{\BankSimStable}{100.0}
\providecommand{\BankWidenInd}{}\renewcommand{\BankWidenInd}{1.25}
\providecommand{\BankWidenDir}{}\renewcommand{\BankWidenDir}{1.01}
\providecommand{\BankSizeIndDelta}{}\renewcommand{\BankSizeIndDelta}{[$-$0.004, 0.770]}
\providecommand{\BankSizeIndSim}{}\renewcommand{\BankSizeIndSim}{[0.071, 1.043]}


This section uses the data analysed by \citet{Kripfganz+Sarafidis:2025}: 350
US banks observed over 36 quarters (2006Q1--2014Q4), distributed with the
replication materials of their article. The dependent variable is the ratio
of non-performing loans to total loans (\code{NPL}), a measure of credit
risk. Operational inefficiency (\code{INEFF}) is treated as endogenous with
respect to the idiosyncratic error and instrumented by the ratio of interest
expenses to deposits (\code{INTEREST}); capital adequacy (\code{CAR}), size
(\code{SIZE}), the capital buffer (\code{BUFFER}), profitability
(\code{PROFIT}), asset quality (\code{QUALITY}) and liquidity
(\code{LIQUIDITY}) are exogenous with respect to the idiosyncratic error but
may be correlated with the factors. The row-standardized weights matrix
links each bank to the banks whose debt ratios are most strongly correlated
with its own. The weights file is an unlabelled matrix ordered like the
bank identifiers, so it is passed together with these identifiers:

\begin{CodeInput}
>>> banks = pd.read_stata("v113i06.dta")
>>> W = pd.read_csv("W.csv", header=None).to_numpy()
>>> ids = sorted(banks["ID"].unique())
>>> common = ["CAR", "SIZE", "BUFFER", "PROFIT", "QUALITY", "LIQUIDITY"]
>>> X, IV = ["INEFF"] + common, ["INTEREST"] + common
>>> spec = dict(W_ids=ids, splag=True, tlags=1, iv=IV, iv_lags=1,
...             iv_splags=True, std=True)
>>> res = sd.fit(banks, "NPL", X, "ID", "TIME", W, **spec)
>>> mg = sd.fit(banks, "NPL", X, "ID", "TIME", W, estimator="mg", **spec)
\end{CodeInput}

The specification has 28 instruments for 9 parameters. All results of this
section are produced by \code{replication/bank\_replication.py}.

\subsection{Cross-software check}

Table~\ref{tab:bank_compare} compares the estimates with those printed in
Tables~1 and~5 of \citet{Kripfganz+Sarafidis:2025}. For both the 2SIV and
the mean-group estimator, the coefficients and standard errors agree with
the published ones in every printed digit: the largest absolute difference
is below $10^{-6}$, the precision of the printed numbers. The $J$ statistic
(\BankJ), the factor and idiosyncratic standard deviations, the numbers of
factors, and the long-run direct, indirect and total impacts of Table~3 of
the original article are reproduced to the same precision. A test in the
package (\code{tests/test\_stata\_replication.py}) checks this agreement
whenever the replication materials are available.

\begin{table}[t!]
\centering\scriptsize
\begin{tabular}{lrrrrrrrr}
\toprule
 & \multicolumn{4}{c}{2SIV (Table~1)} & \multicolumn{4}{c}{Mean group (Table~5)} \\
\cmidrule(lr){2-5}\cmidrule(lr){6-9}
 & \multicolumn{2}{c}{Coefficient} & \multicolumn{2}{c}{Std.\ error} & \multicolumn{2}{c}{Coefficient} & \multicolumn{2}{c}{Std.\ error} \\
\cmidrule(lr){2-3}\cmidrule(lr){4-5}\cmidrule(lr){6-7}\cmidrule(lr){8-9}
 & \pkg{spdfiv} & \proglang{Stata} & \pkg{spdfiv} & \proglang{Stata} & \pkg{spdfiv} & \proglang{Stata} & \pkg{spdfiv} & \proglang{Stata} \\
\midrule
\code{W.NPL} & 0.3943 & 0.3943 & 0.0849 & 0.0849 & 0.0316 & 0.0316 & 0.0511 & 0.0511 \\
\code{L1.NPL} & 0.2899 & 0.2899 & 0.0544 & 0.0544 & 0.3005 & 0.3005 & 0.0149 & 0.0149 \\
\code{INEFF} & 0.4474 & 0.4474 & 0.1046 & 0.1046 & 0.7588 & 0.7588 & 0.1584 & 0.1584 \\
\code{CAR} & 0.0305 & 0.0305 & 0.0058 & 0.0058 & 0.2181 & 0.2181 & 0.0263 & 0.0263 \\
\code{SIZE} & 0.2226 & 0.2226 & 0.0942 & 0.0942 & 2.0040 & 2.0040 & 0.3385 & 0.3385 \\
\code{BUFFER} & $-$0.0545 & $-$0.0545 & 0.0119 & 0.0119 & $-$0.3764 & $-$0.3764 & 0.0420 & 0.0420 \\
\code{PROFIT} & $-$0.0053 & $-$0.0053 & 0.0018 & 0.0018 & $-$0.0180 & $-$0.0180 & 0.0059 & 0.0059 \\
\code{QUALITY} & 0.1830 & 0.1830 & 0.0308 & 0.0308 & 0.2873 & 0.2873 & 0.1387 & 0.1387 \\
\code{LIQUIDITY} & 2.4524 & 2.4524 & 0.2696 & 0.2696 & 6.3302 & 6.3302 & 0.5059 & 0.5059 \\
\midrule
$J$ statistic & 18.825 & 18.825 & & & & & & \\
$\sigma_f$ & 0.6416 & 0.6416 & & & & & & \\
$\sigma_e$ & 0.9038 & 0.9038 & & & & & & \\
\bottomrule
\end{tabular}

\caption{Bank data: estimates from \pkg{spdfiv} and from \code{spxtivdfreg}
as published in Tables~1 and~5 of \citet{Kripfganz+Sarafidis:2025}.}
\label{tab:bank_compare}
\end{table}

Reaching this agreement required five implementation details of
\code{spxtivdfreg} that are not stated in the equations of the article and
that we identified from the source code of the command, distributed under
the GPL with the article. (i) With option \code{std}, variables are
standardized across banks \emph{within each period} before the factors are
extracted. (ii) The eigenvalue-ratio criterion chooses between one and
\code{factmax} factors, without the mock eigenvalue of
\citet{Ahn+Horenstein:2013}. (iii) All lag blocks of a set of instruments
use the largest number of factors selected for any of them. (iv) The
variance of the 2SIV estimator and the $J$ statistic use the weight matrix
computed from the first-stage residuals, as in Equation~(11) of the article.
(v) The instruments are demeaned before, but not after, the factors are
projected out. Point (v) leaves $\hat{\bm Z}_i^\top\bm C_i$ and
$\hat{\bm Z}_i^\top\bm y_i$ unchanged, because the regressors are demeaned,
but it changes the IV weight $\hat{\bm Z}_i^\top\hat{\bm Z}_i$; its effect on
the pooled estimator is of the order $10^{-4}$, whereas for the mean-group
estimator, which uses a separate weight for each bank, it changes the
spatial coefficient from 0.087 to 0.032. Points (i)--(v) are the defaults
of \pkg{spdfiv}; points (i), (ii) and (iv) can be changed through arguments
of \fct{fit} (Table~\ref{tab:options}). Without them, the selected numbers
of factors differ (one instead of two in the instruments, none instead of
one in the residuals), and the coefficient on \code{LIQUIDITY} is 1.50
instead of 2.45. Making such details explicit is, in our view, a useful
by-product of an independent implementation: they matter for the results
but cannot be inferred from the published equations.

\subsection{Long-run impacts}

Table~\ref{tab:bank_impacts} reports the long-run impacts with Delta-method
and simulation intervals (10{,}000 draws, all of which satisfy the stability
condition). The two methods agree for the direct impacts: the simulation
intervals are on average as wide as the Delta intervals, with a median width
ratio of \BankWidenDir. They disagree for the indirect impacts, which
depend on the multiplier $1/(1-\rho_1-\psi_0)$ with $\hat\rho_1+\hat\psi_0=0.68$.
The simulation intervals are skewed to the right, and are \BankWidenInd{}
times as wide as the Delta intervals in the median. The disagreement changes
a conclusion: the Delta interval for the indirect impact of \code{SIZE},
\BankSizeIndDelta, contains zero, whereas the simulation interval,
\BankSizeIndSim, does not. Given the Monte Carlo evidence of
Section~\ref{sec:mc}, we read this as showing that the uncertainty about
the spillovers is asymmetric, not as establishing that the simulation
interval is the more accurate one in this sample.

\begin{table}[t!]
\centering\scriptsize
\begin{tabular}{llrrrrr}
\toprule
 & & & \multicolumn{2}{c}{Delta method} & \multicolumn{2}{c}{Simulation} \\
\cmidrule(lr){4-5}\cmidrule(lr){6-7}
Impact & Variable & Estimate & s.e. & 95\% interval & s.e. & 95\% interval \\
\midrule
direct & \code{INEFF} & 0.647 & 0.159 & [0.335, 0.959] & 0.163 & [0.339, 0.979] \\
 & \code{CAR} & 0.044 & 0.009 & [0.026, 0.062] & 0.009 & [0.027, 0.063] \\
 & \code{SIZE} & 0.322 & 0.142 & [0.044, 0.600] & 0.143 & [0.056, 0.614] \\
 & \code{BUFFER} & $-$0.079 & 0.018 & [$-$0.115, $-$0.043] & 0.018 & [$-$0.117, $-$0.044] \\
 & \code{PROFIT} & $-$0.008 & 0.002 & [$-$0.012, $-$0.003] & 0.002 & [$-$0.012, $-$0.003] \\
 & \code{QUALITY} & 0.265 & 0.047 & [0.173, 0.356] & 0.047 & [0.179, 0.363] \\
 & \code{LIQUIDITY} & 3.547 & 0.445 & [2.674, 4.420] & 0.462 & [2.736, 4.513] \\
\addlinespace
indirect & \code{INEFF} & 0.769 & 0.335 & [0.112, 1.427] & 0.588 & [0.294, 1.909] \\
 & \code{CAR} & 0.052 & 0.024 & [0.006, 0.099] & 0.042 & [0.020, 0.138] \\
 & \code{SIZE} & 0.383 & 0.198 & [$-$0.004, 0.770] & 0.285 & [0.071, 1.043] \\
 & \code{BUFFER} & $-$0.094 & 0.043 & [$-$0.178, $-$0.010] & 0.074 & [$-$0.244, $-$0.035] \\
 & \code{PROFIT} & $-$0.009 & 0.005 & [$-$0.018, 0.000] & 0.007 & [$-$0.024, $-$0.002] \\
 & \code{QUALITY} & 0.315 & 0.141 & [0.039, 0.591] & 0.244 & [0.121, 0.820] \\
 & \code{LIQUIDITY} & 4.218 & 1.742 & [0.803, 7.633] & 3.100 & [1.801, 10.548] \\
\addlinespace
total & \code{INEFF} & 1.417 & 0.427 & [0.579, 2.254] & 0.659 & [0.724, 2.731] \\
 & \code{CAR} & 0.097 & 0.029 & [0.039, 0.154] & 0.046 & [0.052, 0.191] \\
 & \code{SIZE} & 0.705 & 0.310 & [0.097, 1.312] & 0.380 & [0.140, 1.529] \\
 & \code{BUFFER} & $-$0.173 & 0.054 & [$-$0.279, $-$0.066] & 0.083 & [$-$0.343, $-$0.089] \\
 & \code{PROFIT} & $-$0.017 & 0.006 & [$-$0.029, $-$0.004] & 0.008 & [$-$0.034, $-$0.006] \\
 & \code{QUALITY} & 0.580 & 0.167 & [0.252, 0.907] & 0.265 & [0.334, 1.136] \\
 & \code{LIQUIDITY} & 7.765 & 1.904 & [4.034, 11.496] & 3.262 & [5.021, 14.366] \\
\bottomrule
\end{tabular}

\caption{Bank data: long-run average impacts with 95\% intervals from the
Delta method and from Krinsky--Robb simulation (10{,}000 draws).}
\label{tab:bank_impacts}
\end{table}

\subsection{What drives the mean-group estimates?}

\citet{Kripfganz+Sarafidis:2025} found that the MG estimates of some slopes
differ from the 2SIV estimates by an order of magnitude, although the $J$
test does not reject the homogeneous model, and that the spatial spillover
disappears under heterogeneous slopes; they suggested that a few outlying
banks might be responsible. The diagnostics of Section~\ref{sec:diag}
allow this conjecture to be examined.

\begin{table}[t!]
\centering\small
\begin{tabular}{rlrrrr}
\toprule
Bank & Coefficient & $\hat\theta_i$ & Modified $z$ & Change in MG (s.e.) & cond$(\bm Z_i^\top\bm Z_i)$ \\
\midrule
147 & \code{BUFFER} & 3.618 & 7.58 & 0.272 & $1.7\times 10^{6}$ \\
26 & \code{SIZE} & 32.324 & 9.14 & 0.257 & $1.7\times 10^{6}$ \\
112 & \code{SIZE} & 30.991 & 8.76 & 0.245 & $3.7\times 10^{6}$ \\
296 & \code{SIZE} & 30.810 & 8.70 & 0.244 & $2.5\times 10^{7}$ \\
276 & \code{BUFFER} & 2.726 & 5.82 & 0.212 & $1.0\times 10^{7}$ \\
199 & \code{BUFFER} & $-$3.368 & $-$6.22 & $-$0.204 & $1.9\times 10^{7}$ \\
79 & \code{SIZE} & 25.680 & 7.22 & 0.200 & $8.6\times 10^{5}$ \\
16 & \code{BUFFER} & $-$3.207 & $-$5.90 & $-$0.193 & $1.2\times 10^{7}$ \\
\bottomrule
\end{tabular}

\caption{Bank data: the eight banks with the largest influence on the MG
estimates of the \code{SIZE} and \code{BUFFER} slopes. ``Change in MG'' is
the change in the MG estimate from deleting the bank, in units of its
standard error.}
\label{tab:bank_influence}
\end{table}

\begin{table}[t!]
\centering\scriptsize
\begin{tabular}{lcccc}
\toprule
 & MG & Trimmed (10\%) & Median & MG excl.\ flagged \\
\midrule
\code{W.NPL} & 0.032 & 0.020 & 0.034 & 0.023 \\
\footnotesize  & \footnotesize [$-$0.069, 0.132] & \footnotesize [$-$0.063, 0.102] & \footnotesize [$-$0.031, 0.117] & \footnotesize [$-$0.065, 0.109] \\
\code{L1.NPL} & 0.301 & 0.301 & 0.301 & 0.303 \\
\footnotesize  & \footnotesize [0.271, 0.330] & \footnotesize [0.272, 0.334] & \footnotesize [0.268, 0.345] & \footnotesize [0.271, 0.334] \\
\code{INEFF} & 0.759 & 0.507 & 0.313 & 0.191 \\
\footnotesize  & \footnotesize [0.448, 1.069] & \footnotesize [0.271, 0.763] & \footnotesize [0.099, 0.572] & \footnotesize [$-$0.001, 0.403] \\
\code{CAR} & 0.218 & 0.184 & 0.123 & 0.133 \\
\footnotesize  & \footnotesize [0.167, 0.270] & \footnotesize [0.143, 0.225] & \footnotesize [0.094, 0.156] & \footnotesize [0.094, 0.171] \\
\code{SIZE} & 2.004 & 1.282 & 0.807 & 0.707 \\
\footnotesize  & \footnotesize [1.341, 2.668] & \footnotesize [0.796, 1.750] & \footnotesize [0.232, 1.198] & \footnotesize [0.300, 1.123] \\
\code{BUFFER} & $-$0.376 & $-$0.321 & $-$0.220 & $-$0.234 \\
\footnotesize  & \footnotesize [$-$0.459, $-$0.294] & \footnotesize [$-$0.389, $-$0.254] & \footnotesize [$-$0.281, $-$0.167] & \footnotesize [$-$0.301, $-$0.171] \\
\code{PROFIT} & $-$0.018 & $-$0.016 & $-$0.008 & $-$0.020 \\
\footnotesize  & \footnotesize [$-$0.030, $-$0.006] & \footnotesize [$-$0.025, $-$0.007] & \footnotesize [$-$0.018, $-$0.001] & \footnotesize [$-$0.030, $-$0.011] \\
\code{QUALITY} & 0.287 & 0.179 & 0.081 & 0.159 \\
\footnotesize  & \footnotesize [0.015, 0.559] & \footnotesize [0.021, 0.331] & \footnotesize [$-$0.055, 0.214] & \footnotesize [0.011, 0.304] \\
\code{LIQUIDITY} & 6.330 & 4.921 & 3.200 & 3.383 \\
\footnotesize  & \footnotesize [5.339, 7.322] & \footnotesize [4.075, 5.772] & \footnotesize [2.573, 4.049] & \footnotesize [2.691, 4.033] \\
\midrule
Units & 350 & 350 & 350 & 255 \\
\bottomrule
\end{tabular}

\caption{Bank data: MG estimates and robust summaries of the bank-specific
estimates, with 95\% intervals (normal for MG, bootstrap percentile with 999
resamples of banks for the others).}
\label{tab:bank_mg}
\end{table}

Table~\ref{tab:bank_influence} shows that single banks move the MG estimate
of \code{SIZE} or \code{BUFFER} by about 0.19 to 0.27 standard errors. Their
instrument matrices are far from singular (condition numbers below
$10^{8}$), so these are not numerical artefacts. The flags mark
\BankNflagSize{} banks for \code{SIZE}, \BankNflagBuffer{} for
\code{BUFFER}, and \BankNflag{} of the \BankNunits{} banks for at least one
covariate slope. The distribution of the bank-specific estimates is
therefore heavy tailed rather than contaminated by a handful of outliers.

Table~\ref{tab:bank_mg} gives the robust summaries. For \code{SIZE} the
median, \BankMedSize, closes \BankGapCutSize\% of the gap between the MG
estimate (\BankMGsize) and the 2SIV estimate (\BankTwoSize); for
\code{BUFFER} the median closes \BankGapCutBuffer\% of the gap
(\BankMedBuffer{} against \BankMGbuffer{} and \BankTwoBuffer). Excluding the
flagged banks gives similar values (\BankExclSize{} and \BankExclBuffer).
The spatial coefficient behaves differently: the median (\BankMedPsi) and
the MG estimate without the flagged banks (\BankExclPsi) are as small as the
MG estimate itself (\BankMGpsi) and far below the 2SIV estimate
(\BankTwoPsi). Five banks have nearly singular instrument matrices, because
\code{QUALITY} does not vary over time for them, but removing them leaves
the MG estimates practically unchanged (\code{SIZE}: \BankExclIllSize).

The conjecture is thus partly supported. Extreme bank-level estimates
account for about two thirds of the difference in the \code{SIZE} slope and
for about half of it in the \code{BUFFER} slope, but not for the
disappearance of the spatial spillover, which is a feature of the bulk of
the distribution. The MG estimate of the spatial coefficient is also
sensitive to the construction of the instruments: without double
defactoring it is \BankMGpsiNoDD{} instead of \BankMGpsi{}, and point (v)
above moves it as much. One explanation consistent with this evidence,
which we do not pursue here, is the finite-sample behaviour of bank-level
IV estimation with 28 instruments and 35 observations per bank, a ratio at
which IV estimates are known to be pulled towards least squares.

\subsection{Staggered entry}

Finally, the bank panel is made unbalanced artificially, so that the answer
is known: for a random \BankNlate{} of the 350 banks the first 1 to 12
quarters are deleted. Table~\ref{tab:bank_staggered} compares three
estimates: the full balanced panel, the experimental unbalanced estimator
using all \BankNobsStag{} remaining observations, and the balanced
alternative that drops the late banks (\BankNobsDrop{} observations).

\begin{table}[t!]
\centering\scriptsize
\begin{tabular}{lrrrrrr}
\toprule
 & \multicolumn{2}{c}{Full panel} & \multicolumn{2}{c}{Staggered, all banks} & \multicolumn{2}{c}{Late banks dropped} \\
\cmidrule(lr){2-3}\cmidrule(lr){4-5}\cmidrule(lr){6-7}
 & Coef. & s.e. & Coef. & s.e. & Coef. & s.e. \\
\midrule
\code{W.NPL} & 0.394 & 0.085 & 0.352 & 0.106 & 0.367 & 0.105 \\
\code{L1.NPL} & 0.290 & 0.054 & 0.305 & 0.054 & 0.243 & 0.058 \\
\code{INEFF} & 0.447 & 0.105 & 0.435 & 0.101 & 0.456 & 0.117 \\
\code{CAR} & 0.031 & 0.006 & 0.032 & 0.006 & 0.030 & 0.006 \\
\code{SIZE} & 0.223 & 0.094 & 0.199 & 0.105 & 0.330 & 0.124 \\
\code{BUFFER} & $-$0.055 & 0.012 & $-$0.051 & 0.013 & $-$0.055 & 0.013 \\
\code{PROFIT} & $-$0.005 & 0.002 & $-$0.004 & 0.002 & $-$0.006 & 0.002 \\
\code{QUALITY} & 0.183 & 0.031 & 0.173 & 0.031 & 0.182 & 0.029 \\
\code{LIQUIDITY} & 2.452 & 0.270 & 2.525 & 0.282 & 2.468 & 0.325 \\
\midrule
Observations & \multicolumn{2}{c}{12{,}250} & \multicolumn{2}{c}{11{,}602} & \multicolumn{2}{c}{8{,}575} \\
Banks & \multicolumn{2}{c}{350} & \multicolumn{2}{c}{350} & \multicolumn{2}{c}{245} \\
\bottomrule
\end{tabular}

\caption{Bank data with artificial staggered entry: 2SIV estimates.}
\label{tab:bank_staggered}
\end{table}

The unbalanced estimates stay within \BankMaxStagDiffSE{} standard errors
of the full-panel estimates. Dropping the late banks moves the estimates by
up to \BankMaxDropDiffSE{} standard errors, because it discards 30\% of the
network as well as the observations. In this application the experimental
option thus recovers more of the full-sample information than the balanced
alternative available in existing software.


%% ============================================================================
\section{Computing time} \label{sec:timing}
%% ============================================================================

Table~\ref{tab:timing} reports the best of three wall-clock times, in
seconds, for the 2SIV and MGIV estimators and for long-run impacts with
1{,}000 simulation draws, on the platform described in the computational
details. Estimation with a thousand units and fifty periods takes less than half a second; the impact computation is dominated by the Schur decomposition of $\bm W$, which is computed once.

\begin{table}[t!]
\centering\small
\begin{tabular}{rrrrr}
\toprule
$N$ & $T$ & 2SIV & Impacts (1000 draws) & MGIV \\
\midrule
100 & 50 & 0.03 & 0.02 & 0.04 \\
250 & 50 & 0.06 & 0.06 & 0.09 \\
500 & 50 & 0.11 & 0.26 & 0.19 \\
1000 & 50 & 0.33 & 1.57 & 0.48 \\
100 & 200 & 0.05 & 0.01 & 0.06 \\
500 & 200 & 0.23 & 0.25 & 0.32 \\
\bottomrule
\end{tabular}

\caption{Computing time in seconds (best of three runs).}
\label{tab:timing}
\end{table}

%% ============================================================================
\section{Summary and discussion} \label{sec:summary}
%% ============================================================================

\pkg{spdfiv} makes the defactored IV approach to spatial dynamic panels
with interactive effects available as open-source software. It reproduces
the defining equations of the 2SIV estimator to machine precision. On the data of
\citet{Kripfganz+Sarafidis:2025} it reproduces the published 2SIV
and mean-group results of \code{spxtivdfreg} in every printed digit, and
the comparison made explicit five implementation details of that command. The package adds safeguards and
diagnostics that matter in applied work: label-based
alignment of the weights matrix, exact and fast computation of impacts with
a choice between Delta-method and simulation inference, tools that show
which units drive mean-group estimates, and an experimental treatment of
unbalanced panels.

Several limitations remain. The asymptotic theory for unbalanced panels and
for the robust MG summaries has not been developed; the corresponding
options are diagnostics and are labelled as such. The package does not
implement a test of remaining cross-sectional dependence, because the
standard CD test is invalid when factors have been estimated from the same
data \citep{Juodis+Reese:2022}. Planned extensions include an interface for
\pkg{PySAL} weights objects, a test of slope homogeneity, and an \proglang{R}
port.

%% ============================================================================
\section*{Computational details}
%% ============================================================================

The results in this paper were obtained using \proglang{Python}~3.13.17
with \pkg{spdfiv}~0.1.0, \pkg{NumPy}~2.5.3, \pkg{SciPy}~1.18.1, and
\pkg{pandas}~3.0.5, on Linux with two CPU cores. The package and the
replication materials are available from \url{https://github.com/abdulameert-boop/spdfiv}.

\section*{Acknowledgments}

\todo{Acknowledgments.} Parts of the software and of the manuscript were
drafted with the assistance of an AI system (Claude, Anthropic); the authors
have checked and tested, and take full responsibility for all code and text.
\todo{Adapt this disclosure to the journal's current policy on AI tools.}

\bibliography{refs}

\end{document}
